\par\noindent\textbf{(1)} 对频率ω的振子定义a\_ω=√(mω/(2ℏ))x̂+i p̂/√(2ℏmω)。频率骤变只改变Hamiltonian系数，x̂与p̂作为同一物理算符在t=0连续；因此要用旧a\_i,a\_i†展开新的a\_f，而不是把旧基态误认为新基态。
\par\noindent\textbf{(2)} 代入x̂=√(ℏ/(2mωi))(a\_i+a\_i†)、p̂=−i√(ℏmωi/2)(a\_i−a\_i†)，得a\_f=α a\_i+β a\_i†，α=(√(ωf/ωi)+√(ωi/ωf))/2、β=(√(ωf/ωi)−√(ωi/ωf))/2，满足α²−β²=1以保持玻色对易。
\par\noindent\textbf{(3)} 旧基态满足a\_i|0\_i>=0，因此新激发数⟨a\_f†a\_f⟩=β²=[ωf/ωi+ωi/ωf−2]/4≥0。新的平均能量⟨H\_f⟩=ℏωf(β²+1/2)=ℏ(ωf²+ωi²)/(4ωi)，比新基态能高，表现为压缩态激发。
\par\noindent\textbf{(4)} ωf→ωi时α→1、β→0，无额外量子数且平均能量回到ℏωi/2。外部骤改弹簧刚度对系统做功W=⟨H\_f⟩−ℏωi/2=ℏ(ωf²−ωi²)/(4ωi)，可正可负；额外激发能不等于全部外功，因为新零点能本身也变化。
